% SEGMENT OLP-0485-S01
% Part: applied-modal-logics
% Chapter: epistemic-logic
% Section: relational-models

\documentclass[../../../include/open-logic-section]{subfiles}

\begin{document}

\olfileid{aml}{el}{rel}

\olsection{தொடர்புமுறை மாதிரிகள்}

அறிவுநிலைத் தருக்கங்களின் அடிப்படைப் பொருண்மையியல் கருத்து, வழக்கமான
வாய்ப்புநிலைத் தருக்கங்களின் கருத்தையே ஒத்தது. தொடர்புமுறை மாதிரிகள் இன்னமும்
உலகங்களின் ஒரு கணத்தையும், எந்த உலகங்களில் எந்த !!{propositional variable}s
``மெய்'' என்று கணக்கிடப்படுகின்றன என்பதைத் தீர்மானிக்கும் மதிப்பளிப்பையும்
கொண்டிருக்கின்றன. ஒரே முகவரை மட்டுமே கையாள்கிறோம் எனில், வழக்கம்போல ஒரே ஓர்
அணுகுறவு இருக்கும். ஆனால் நம்மிடம் பன்முகவர் அறிவுநிலைத் தருக்கம் இருந்தால்,
அந்த ஒற்றை அணுகுறவு அணுகுறவுகளின் ஒரு கணமாகிறது; நமது முகவர்-குறிகளின்
கணம்~$G$ இலுள்ள ஒவ்வொரு~$a$ க்கும் ஒன்று இருக்கும்.

ஒரு \emph{தொடர்புமுறை மாதிரி}, உலகங்களின் கணத்தையும்---ஒவ்வொரு
முகவருக்கும் ஒன்று வீதம் ஈருறுப்பு அணுகுறவுகளால் அவை தொடர்புபடுத்தப்படும்---
எந்த உலகங்களில் எந்த !!{propositional variable}s மெய் என்பதைத் தீர்மானிக்கும்
ஒரு மதிப்பளிப்பையும் கொண்டுள்ளது.

\begin{defn}
  பன்முகவர் அறிவுநிலை மொழிக்கான ஒரு \emph{மாதிரி} என்பது
  $\mModel{M} = \tuple{W, R, V}$ என்ற மும்மை; இதில்
  \begin{enumerate}
  \item $W$ என்பது ``உலகங்களின்'' காலியல்லாத கணம்,
  \item ஒவ்வொரு $a \in G$ க்கும் ${R}_a$ என்பது~$W$ மீதான ஓர் ஈருறுப்பு
    அணுகுறவு, மேலும்
  \item ஒவ்வொரு !!{propositional variable}~$p$ க்கும் சாத்தியமான உலகங்களின்
    ஒரு கணம் $V(p)$ ஐ மதிப்பளிக்கும் சார்பு $V$.
  \end{enumerate}
  $R_a ww'$ ஆகும்போது, $w'$ ஆனது~$w$ இலிருந்து \emph{a ஆல்
  அணுகத்தக்கது} என்கிறோம். $w \in V(p)$ ஆகும்போது, $p$ ஆனது~$w$ இல்
  \emph{மெய்} என்கிறோம்.
\end{defn}

கூடுதலான அணுகுறவுகள் சேர்க்கப்பட்டுள்ளதைத் தவிர, இயங்குமுறை இயல்பான
வாய்ப்புநிலைத் தருக்கத்தின் இயங்குமுறையைப் போன்றதே. கொடுக்கப்பட்ட ஒரு
முகவருக்கு, அவரது அணுகுறவு அவரது தகவல் நிலைகள் குறித்த ஏதோவொன்றைக்
குறிப்பதாகப் பொதுவாகப் பொருள்கொள்வோம். எடுத்துக்காட்டாக, $R_a ww'$ என்பது
$w'$ ஆனது~$w$ இல்~$a$ கொண்டுள்ள தகவலுடன் ஒத்திசைவானது என்பதை
வெளிப்படுத்துவதாக அடிக்கடி எடுத்துக்கொள்கிறோம். வேறுவிதமாகச் சொன்னால்,
$w$ இல் அவரால் உலகம்~$w$ க்கும் உலகம்~$w'$ க்கும் இடையில் வேறுபாடு
காண முடியாது.

\end{document}
